\documentclass[11pt]{article}
\usepackage{amsmath,amssymb,amsthm}
\newtheorem{theorem}{Theorem}
\newtheorem{lemma}{Lemma}
\title{GST true constant-family control: countable-field descent of every EXISTING rational algebraic Betti class and the exact source-birth boundary}
\date{12 November 2026}
\begin{document}
\maketitle

\section{The user's concrete universal request}
Fix an arbitrary smooth projective complex ORIGINAL X,
codimension p and a given rational (p,p) Hodge class
alpha. The owner wants to construct an actual original-X
cycle Z_alpha with canonical Betti class alpha,
using arbitrarily many GST original U,J,M,[J,M],
Navigation, N-cohomology and cardinal-indexed worlds.
The proposed new mechanism is to MANUFACTURE
a genuine constant family rather than assume
uncountably many distinct original-X deformations
already carry target-matching cycles.

The exact distinction to resolve is:
a constant family X times B CAN transport and
SPECIALIZE real algebraic cycles along a parameter,
but can it produce one from a prescribed Hodge class
with no algebraic source input?

\section{Unconditional positive theorem: every algebraic cohomology class has a countable-field witness}
\begin{theorem}[Actual original-X Betti-source descent to countable constants]
For EVERY smooth projective complex X there exists
a COUNTABLE algebraically closed subfield
k of C and a smooth projective k-variety X_k
with X_k times_k C isomorphic to X, such that
for EVERY p and EVERY actual rational algebraic
cycle Z in CH^p(X)_Q there is an actual rational
cycle Z_k in CH^p(X_k)_Q satisfying
\[
\boxed{
\mathrm{cl}^{can}_{X}
 \bigl((Z_k)_{\mathbb C}\bigr)
 =
\mathrm{cl}^{can}_{X}(Z)
\quad\text{in }H^{2p}(X,\mathbb Q).
}
\tag{1}
\]
The stronger assertion Z_k and Z must agree
in the Chow GROUP is NOT asserted;
only their genuine canonical rational Betti classes
are required and proved.
\end{theorem}

\begin{proof}
Embed X projectively over C. Its finitely many
defining equations have finitely many complex
coefficients c1,...,cn. Put
k=algebraic closure inside C of Q(c1,...,cn).
This field is COUNTABLE. Smoothness and
projectivity descend to a model X_k by
faithfully flat base change.

Choose an actual rational Chow cycle
Z=(1/m)(P-Q) with finite effective algebraic
cycle sums P,Q and positive integer m.
Their finitely many defining equations,
component multiplicities and projective
Hilbert data use finitely many more complex
coefficients. Take a finitely generated
integral k-subalgebra A of C containing
these coefficients, and set S=Spec A.
Spread the corresponding components to
actual closed subschemes of X_k times S.
By generic flatness, shrink to a nonempty
Zariski open U subset S where the components
are flat with fixed relative Hilbert data.
The original C-point of S lies in U after
the usual suitable model/shrink construction.

The base U is integral, hence U(C) is
connected in the classical complex topology.
The relative projective cycles have
locally constant canonical rational Betti
classes in the true constant smooth
family X(C) times U(C). This follows
from their global relative cohomology
classes and homotopy invariance of
fiber restrictions; no Hodge conjecture
or algebraic source-birth theorem is used.

Because k is algebraically closed and U
is a nonempty finite-type k-scheme, U has
a k-rational closed point u
(by Nullstellensatz). Evaluate the
actual flat relative cycles at u.
Their difference divided by m gives
a genuine cycle Z_k on the ORIGINAL
k-model X_k. Since the Betti class
was locally constant and U(C) connected,
its complex base change has EXACTLY
the same canonical rational Betti class
as the original Z. This proves (1).
\end{proof}

\section{New EXACT infinite-control consequence}
For the fixed k-model X_k, the set of
finite projective equations with k-coefficients
is COUNTABLE. Thus actual finite rational
Chow cycles over k are COUNTABLE as a set.
The theorem shows that the entire rational
canonical Betti Chow IMAGE of X_C is
generated by this COUNTABLE explicit
universe of actual k-defined cycles:
\[
\boxed{
\operatorname{im}(\mathrm{cl}_X)
=
\bigl\{
 \mathrm{cl}_X(Z):
 Z\in CH^p(X_k)_{\mathbb Q}
 \bigr\}.
}\tag{2}
\]
For each p, choose an enumeration of ALL
finite k-defined algebraic cycles
Z0,Z1,Z2,... with rational coefficients.
Then, for every algebraic rational class
beta in the image, SOME finite index n
has cl(Z_n)=beta.

This is a genuine arithmetic/constant-family
normal form, independent of the number of
GST cardinal stages and valid for ALL
smooth projective complex original X.
It supplies countably many ACTUAL,
not synthetic, original-X algebraic
candidate sources, and proves any
EXISTING algebraic class is hit by one.
It does NOT prove that a prescribed
Hodge class alpha is hit.

\section{Why native infinite enumeration has no automatic termination}
The exact missing statement would be
\[
\forall\alpha\in W_p(X),\
 \exists n\in\mathbb N:
 \mathrm{cl}(Z_n)=\alpha.
\]
This is exactly the rational Hodge
conjecture after the countable-field
descent theorem, NOT an independently
proved step. A transfinite/aleph-indexed
search through the countable k-cycle list
cannot guarantee that a match exists,
regardless of the speed, U/J/M curvature,
navigation gating or source orbit rank.

If a native GST theorem provided an
independent target-sensitive termination
certificate (with a real algebraic cycle
construction from alpha and verified
canonical Betti class), that would be
a complete proof. The present theorem
makes a concrete smaller arithmetic
search domain but does not manufacture
such a certificate.

\section{Comparison with algebraic specialization across constant families}
The construction in the proof makes a
TRUE constant smooth projective family
X_k times U and genuine relative
algebraic cycles already present
over its generic/complex parameter.
Its geometric specializations
over k-points preserve the cohomology
class. This is a real positive
constant-family descent theorem.

Conversely, the existence of an actual
Z on a generic fiber of X times U with
canonical class alpha is already
an EXISTENCE input in that generic fiber.
Restricting to any fiber by its actual
refined Gysin map returns an original-X
cycle; varying the topology or cardinal
number of OTHER parameter points
does not create Z when it was absent.

\section{Proof/evidence status}
The field-countability and specialization
proof is a complete classical algebraic
geometry argument at hand-proof level.
No new Lean assertion is claimed to
construct the Hilbert/Chow parameter
spaces automatically. The theorem has
not resolved the universal given-alpha
original-X source existence arrow.
\end{document}